\chapter{Systems-and-control reduction: balanced truncation, Krylov and SOAR}\label{ch:bt}

\begin{tcolorbox}[colback=gray!5,colframe=gray!50,title={In this chapter}]
A damped cantilever is going to be part of a larger system, and all the larger system will ever ask of it is one thing: apply a force at the tip, read the tip displacement, at many frequencies. Chapter~\ref{ch:pod} reduced a model by looking at its solutions; here we look only at that input--output map. Three families do it. Balanced truncation keeps the states that are both easy to excite and easy to see, and comes with an error bound known in advance. Krylov methods match the response exactly at one frequency and need no Lyapunov solve. SOAR does the same without leaving the second-order form of structural dynamics. We run them on the beam against an independent full-order sweep and measure the bound, the accuracy at equal size, how local moment matching is, a scale problem in the SOAR code, and stability and passivity.\par\smallskip
After it you should be able to choose among the three for a given port, check the bound and the stability of a reduced model, and know that moment matching is accurate only near its expansion point.
\end{tcolorbox}

\section{The problem}\label{sec:bt-problem}

Imagine the beam as one component of a bigger machine. The machine pushes on its tip and asks how far the tip moves, at four hundred frequencies, every time a design changes. The direct route is a loop over frequencies, a complex linear solve at each. It is correct and it is the reference in this chapter, but the cost is paid again for every frequency and every design, and the solve does the work of finding all the internal motion of the beam although only one number leaves it. A reduced model should hand over just that number.

A structural model with damping is $\bm M\ddot{\bm q}+\bm C\dot{\bm q}+\bm K\bm q=\bm b\,u(t)$ with output $y=\bm c^T\bm q$: one force $u$ acting at a chosen degree of freedom, one displacement $y$ observed. All that matters to a user of such a port is its transfer function
\begin{equation}\label{eq:bt-tf}
H(s)=\bm c^T\bigl(s^2\bm M+s\bm C+\bm K\bigr)^{-1}\bm b,\qquad s=i\omega,
\end{equation}
and the task is to find a model of much smaller order whose $H_r(s)$ is close to $H(s)$ over the band of interest. The dimension is not the number of unknowns: the first-order form used by the state-space methods has $2n$ states, $\bm x=[\bm q;\dot{\bm q}]$, and its size is the unit in which the methods below are compared (\emph{state count}; a second-order model of $k$ unknowns has $2k$ states).

\paragraph{The model problem.}
A cantilever of $24$ Euler--Bernoulli elements (\cbtNDof{} free unknowns, \cbtNState{} states) with Rayleigh damping $\bm C=\alpha\bm M+\beta\bm K$, $\alpha=2$, $\beta=10^{-5}$, the model of the package example \path{rom_engine/examples/gallery/rom_04_balanced_truncation.py}. The force and the observed displacement are both at the tip, so $\bm b=\bm c$ (a \emph{collocated} port). The first natural frequency is \cbtFOne{}~Hz. The reference is the full-order frequency sweep of \fea{} (\code{solve_frequency_sweep}), independent of everything in \rom{}, on $400$ frequencies from $0.05\omega_1$ to $8\omega_1$.

\section{Mathematical structure}\label{sec:bt-math}

\begin{sidenote}{Transfer function and state count}
The transfer function $H(s)$ says what comes out at the output for a unit sinusoidal force at frequency $\omega$, with $s=i\omega$. The state-space methods write the second-order equation as a first-order one with state $[\bm q;\dot{\bm q}]$, which has twice the unknowns; to compare methods fairly, sizes are counted in states, and a second-order model with $k$ unknowns has $2k$ states.
\end{sidenote}

\begin{sidenote}{Gramians and Hankel singular values in words}
Think of a model with an input (a force) and an output (a measured displacement). Some states are easy to excite with the input but hardly visible in the output, and others show clearly in the output but are hard to excite. The \emph{controllability Gramian} measures how much input energy is needed to reach a state, and the \emph{observability Gramian} how much output energy a state produces. In balanced coordinates both are equal and diagonal, so every state is as hard to excite as it is to see, and its \emph{Hankel singular value} says how much it matters to the input--output behaviour. Truncation discards the states with the smallest values, and the error bound below is a multiple of their sum. The method requires a \emph{stable} system, one whose free response decays.
\end{sidenote}

\paragraph{Balanced truncation.}
Write the first-order form $\dot{\bm x}=\bm A\bm x+\bm B u$, $y=\bm C_o\bm x$ with $\bm A$ stable. The controllability and observability Gramians solve two Lyapunov equations,
\begin{equation}\label{eq:bt-lyap}
\bm A\bm P+\bm P\bm A^T+\bm B\bm B^T=\bm 0,\qquad \bm A^T\bm Q+\bm Q\bm A+\bm C_o^T\bm C_o=\bm 0 ,
\end{equation}
and measure, respectively, how much a state can be excited by the input and how much it shows at the output. A change of coordinates can make both equal and diagonal, $\bm P=\bm Q=\mathrm{diag}(\sigma_1,\dots,\sigma_{2n})$; the $\sigma_k$ are the Hankel singular values, and the states with small $\sigma_k$ matter least to the input--output map~\cite{moore1981,antoulas2005}. Truncating the balanced model to its first $r$ states keeps the stability of the full model and satisfies the a priori bound~\cite{enns1984,glover1984}
\begin{equation}\label{eq:bt-bound}
\sup_{\omega}\;|H(i\omega)-H_r(i\omega)|\;\le\;2\sum_{k>r}\sigma_k .
\end{equation}
The bound is absolute (in the units of $H$, here m/N), holds for every frequency at once, and is known before any frequency response is evaluated. The method needs two Lyapunov solves of dense $2n\times2n$ systems, so it is a method for models of moderate size or for models that have already been reduced.

\begin{algbox}{Square-root balanced truncation}
\begin{pseudo}
\Require stable $(\bm A,\bm B,\bm C_o)$ of order $N$, target order $r$ or tolerance
\State solve the two Lyapunov equations \eqref{eq:bt-lyap} for $\bm P,\bm Q$;\ Cholesky $\bm P=\bm S\bm S^T$, $\bm Q=\bm R\bm R^T$
\State SVD $\bm R^T\bm S=\bm U\bm\Sigma\bm V^T$;\ Hankel singular values $\sigma_k=\Sigma_{kk}$
\State choose $r$ so that $2\sum_{k>r}\sigma_k\le\mathrm{tol}$ (if not given)
\State $\bm T\gets\bm S\bm V_r\bm\Sigma_r^{-1/2}$,\ $\bm W\gets\bm R\bm U_r\bm\Sigma_r^{-1/2}$
\State $\bm A_r\gets\bm W^T\bm A\bm T$,\ $\bm B_r\gets\bm W^T\bm B$,\ $\bm C_r\gets\bm C_o\bm T$
\Ensure reduced model of order $r$ and the a priori bound \eqref{eq:bt-bound}
\end{pseudo}
\end{algbox}

The cost is two dense Lyapunov solves, $O(N^3)$ in time and $O(N^2)$ in memory, with $N=2n$ states. The method fails when $\bm A$ is not stable, since then there are no Gramians, and when the Gramians are numerically singular or badly conditioned, as for the raw model of \Cref{sec:bt-raw}. The code is \anchor{rom_engine.balanced_truncation.BalancedTruncationROM} and \anchor{rom_engine.balanced_truncation.hankel_singular_values}; \Cref{fig:bt-bt} and \Cref{fig:bt-compare} are the checks.

\begin{sidenote}{Moments and Krylov subspaces}
The transfer function $H(s)$ can be expanded in a power series about a frequency $s_0$, and the coefficients are its \emph{moments}. A \emph{Krylov subspace} is spanned by a vector and its repeated images under a matrix, $\{\bm b,\bm A\bm b,\bm A^2\bm b,\dots\}$, and these images are, up to the factorisation of $s_0\bm E-\bm A$, the moments themselves. Projecting onto an $r$-dimensional Krylov subspace therefore reproduces the first moments of the full model, which makes the reduced model accurate near $s_0$ at the price of a single sparse factorisation. The accuracy is local, however, and there is no global error bound.
\end{sidenote}

\paragraph{Krylov moment matching.}
Expand $H$ about a point $s_0$: $H(s_0+\sigma)=\sum_{j\ge0}m_j\sigma^j$. The moments $m_j$ are the entries of a Krylov sequence built with one sparse factorisation of $s_0\bm E-\bm A$ (the descriptor form $\bm E\dot{\bm x}=\bm A\bm x+\bm B u$ avoids inverting $\bm M$), so an orthonormal basis $\bm V$ of that Krylov subspace, used as a Galerkin projection $\bm A_r=\bm V^T\bm A\bm V$, $\bm E_r=\bm V^T\bm E\bm V$, makes the reduced transfer function match the first moments at $s_0$ exactly~\cite{besselink2013comparison}. A second basis $\bm W$ from the dual system gives the oblique (two-sided, Petrov--Galerkin) projection, which matches about twice as many moments for the same size. Matching is \emph{local}: exact in the limit at $s_0$, and nothing in the construction controls the error elsewhere, and the stability of the full model is not inherited.

\begin{algbox}{Arnoldi basis for moment matching at $s_0$}
\begin{pseudo}
\Require $\bm E,\bm A,\bm B$, expansion point $s_0$, number of vectors $k$, deflation tolerance
\State factor $\bm F\gets s_0\bm E-\bm A$ \emph{once} (sparse)
\State $\bm v\gets\bm F^{-1}\bm B$;\ $\bm V\gets[\bm v/\|\bm v\|]$
\For{$j=2,\dots,k$}
  \State $\bm w\gets\bm F^{-1}\bm E\bm v_{j-1}$ \Comment{next moment direction}
  \State orthogonalise $\bm w$ against the columns of $\bm V$ (modified Gram--Schmidt, twice)
  \If{$\|\bm w\|\le$ tolerance} \textbf{deflate and stop} \Else{} append $\bm w/\|\bm w\|$ to $\bm V$ \EndIf
\EndFor
\State $\bm A_r\gets\bm V^T\bm A\bm V$,\ $\bm E_r\gets\bm V^T\bm E\bm V$,\ $\bm B_r\gets\bm V^T\bm B$ \Comment{one-sided; two-sided adds $\bm W$ from the dual system}
\Ensure basis $\bm V$ (the number of columns returned may be less than $k$), reduced model
\end{pseudo}
\end{algbox}

One sparse factorisation and $k$ back-substitutions are needed, plus $O(Nk^2)$ for the orthogonalisation. It fails if $s_0\bm E-\bm A$ is singular or the tolerance deflates early, so the number of returned vectors has to be checked; away from $s_0$ no error is controlled and stability is not inherited. The code is \anchor{rom_engine.krylov.arnoldi_basis} and \anchor{rom_engine.krylov.KrylovROM}; \Cref{fig:bt-locality} and \Cref{tab:bt-compare} are the checks.

\begin{sidenote}{Second-order and first-order form}
A structural model is naturally of second order, $\bm M\ddot{\bm q}+\bm C\dot{\bm q}+\bm K\bm q=\bm B\bm u$. The control literature uses the first-order form $\dot{\bm x}=\bm A\bm x+\bm B\bm u$, with the state $\bm x=[\bm q;\dot{\bm q}]$, which doubles the number of states and hides the physical structure of the matrices. Methods that work directly with $\bm M$, $\bm C$ and $\bm K$ keep the reduced model a mechanical system with its own mass, damping and stiffness.
\end{sidenote}

\paragraph{SOAR.}
The first-order form doubles the state count. The second-order Arnoldi method of Bai and Su~\cite{bai2005soar} builds a basis of the same moments directly in the pencil $s^2\bm M+s\bm C+\bm K$, with the two-term recurrence $\bm r_j=\bm A_1\bm r_{j-1}+\bm A_2\bm r_{j-2}$, $\bm A_1=-\bm K_0^{-1}(\bm C+2s_0\bm M)$, $\bm A_2=-\bm K_0^{-1}\bm M$, $\bm K_0=\bm K+s_0\bm C+s_0^2\bm M$, and projects $\bm M,\bm C,\bm K$ by congruence. The reduced model is again second-order, with $k$ unknowns and $2k$ states.

\begin{algbox}{SOAR basis in the second-order pencil}
\begin{pseudo}
\Require $\bm M,\bm C,\bm K$, port vector $\bm b$, point $s_0$, size $k$
\State $\bm K_0\gets\bm K+s_0\bm C+s_0^2\bm M$;\ factor once;\ $\bm A_1\gets-\bm K_0^{-1}(\bm C+2s_0\bm M)$,\ $\bm A_2\gets-\bm K_0^{-1}\bm M$ (applied, never formed)
\State $\bm r_0\gets\bm K_0^{-1}\bm b$;\ $\bm r_{-1}\gets\bm 0$
\For{$j=1,\dots,k$}
  \State $\bm r_j\gets\bm A_1\bm r_{j-1}+\bm A_2\bm r_{j-2}$ (for $j=1$ use $\bm r_0$)
  \State orthogonalise $\bm r_j$ against the columns of $\bm\Phi$; append if its norm exceeds the \emph{absolute} tolerance, otherwise deflate
\EndFor
\State $\bm M_r,\bm C_r,\bm K_r\gets\bm\Phi^T\bm M\bm\Phi,\ \bm\Phi^T\bm C\bm\Phi,\ \bm\Phi^T\bm K\bm\Phi$ \Comment{congruence}
\Ensure second-order reduced model with $k$ unknowns ($2k$ states)
\end{pseudo}
\end{algbox}

The cost is one sparse factorisation and $k$ solves, without the doubling of the state. The pitfall is a deflation tolerance that does not match the scale of the data, which the code reports only through the shape of the result (\Cref{sec:bt-soar}); block starting vectors need the TOAR variant, which is not implemented. The code is \anchor{rom_engine.soar.soar_basis} and \anchor{rom_engine.soar.SOARROM}; \Cref{fig:bt-response}, \Cref{tab:bt-compare} and the passivity grid check in \Cref{tab:bt-stab} are the checks.

\begin{sidenote}{Passive}
A structure is passive if it can only dissipate the energy put in through the port, never create it. A reduced model that is not passive can feed energy back when coupled to other components, which makes the combined simulation unstable even if each part looks fine alone.
\end{sidenote}

\paragraph{Passivity.}
For a collocated port the force $u$ and the velocity $y_v=\bm b^T\dot{\bm q}$ have power $uy_v$, and the energy balance $\dot E=-\dot{\bm q}^T\bm C\dot{\bm q}+uy_v\le uy_v$ says the structure only dissipates. Because a congruence $\bm\Phi^T\bm C\bm\Phi$ keeps positive semidefiniteness for \emph{any} full-column-rank basis, every second-order Galerkin projection preserves this inequality; this is the argument in the module docstring of \anchor{rom_engine.passivity.passivity_margin} and it applies to modal truncation, POD (Chapter~\ref{ch:pod}) and SOAR. It does not apply to the first-order projections of balanced truncation and Krylov methods, for which passivity has to be checked. The frequency-domain test for a passive one-port is $\mathrm{Re}\,H_v(i\omega)\ge0$ with $H_v=i\omega H$, or equivalently that the phase of the mobility stays within $\pm90^\circ$; the experiments report $\min_\omega\cos\arg H_v$.

\tikzfigure[0.98\linewidth]{ch09_map}{Map of the systems-and-control methods of this chapter: modal pre-reduction followed by balanced truncation for moderate models, Krylov or SOAR for large sparse ones; every route ends at the same check against data not used to build it.}{fig:bt-mapidea}

\begin{algbox}{Choosing and running a systems-and-control reduction of a structural port}
\begin{pseudo}
\Require port ($\bm b,\bm c$), band, tolerance, full model size
\State check that the damping is not zero \Comment{the Lyapunov equations need a stable $\bm A$}
\If{the model is small, or has been pre-reduced}
  \State pre-reduce by modal truncation if frequencies span many orders of magnitude (\Cref{sec:bt-raw});\ run square-root balanced truncation; choose $r$ from \eqref{eq:bt-bound}
\Else
  \State choose $s_0$ inside the band;\ build Arnoldi or SOAR basis with one sparse factorisation
\EndIf
\State \textbf{check:} pole locations, number of basis vectors returned, a frequency response not used in the construction
\If{structure is passive and the model will be coupled into a larger simulation}
  \State prefer SOAR (second-order congruence) and check the passivity margin
\EndIf
\Ensure reduced port model with its check
\end{pseudo}
\end{algbox}

The cost is that of the chosen branch, a dense $O(N^3)$ solve or one sparse factorisation. The choice only works if the checks are made: moment matching has no error control away from $s_0$, and the first-order projections are not passive by construction. The experiments of \Cref{sec:bt-verif}, summarised in \Cref{tab:bt-compare}, make them.

\section{Implementation in the packages}\label{sec:bt-impl}

The three method families share one entry convention, \code{from_MCK(M, K, B, Cout, C=..., ...)}, and one output convention, \code{frequency_response(omegas)}, so that the same comparison code runs on all of them (\Cref{lst:bt-run}). \Cref{tab:bt-map} lists the code. The package also contains the singular-perturbation variant of balanced truncation, which keeps the static gain exactly, a frequency-weighted variant, and an optimal Hankel-norm variant; these are \emph{not} tested in this book.

\begin{lstlisting}[caption={The three reductions on one collocated port; M, K, C are the free-dof matrices and B the tip load vector.},label={lst:bt-run}]
import numpy as np
from scipy.linalg import eigh
from rom_engine.balanced_truncation import BalancedTruncationROM
from rom_engine.krylov import KrylovROM
from rom_engine.soar import soar_basis, SOARROM

# balanced truncation, after modal pre-reduction to 15 modes
lam, Vf = eigh(K, M); Vp = Vf[:, :15]
Mp, Kp, Cp, Bp = Vp.T @ M @ Vp, Vp.T @ K @ Vp, Vp.T @ C @ Vp, Vp.T @ B
bt = BalancedTruncationROM.from_MCK(Mp, Kp, Bp, Bp, C=Cp, r=8)
print(bt.h_infinity_error_bound(), bt.is_stable())

# Krylov about s0 = 0 on the FULL model: 8 states, one-sided and two-sided
k1 = KrylovROM.from_MCK(M, K, B, B, C=C, s0=0.0, k=8)
k2 = KrylovROM.from_MCK(M, K, B, B, C=C, s0=0.0, k=4, two_sided=True)

# SOAR: 4 unknowns = 8 states; load and time rescaled first (see text)
w1 = np.sqrt(lam[0])
V = soar_basis(w1**2 * M, K, B / np.linalg.norm(np.linalg.solve(K, B)), 4, C=w1 * C)
soar = SOARROM(M, C, K, B.reshape(-1, 1), B.reshape(1, -1), V, 0.0)

H = [m.frequency_response(omegas) for m in (bt, k1, k2, soar)]
\end{lstlisting}

\paragraph{Try it.} \Cref{lst:bt-try} sets the loop against two eight-state models of the same beam: the loop first, then balanced truncation (after modal pre-reduction) and Krylov at $s_0=0$. The beam is the one of the example script; the last two lines of output are the ones to read.

\begin{lstlisting}[caption={The naive frequency loop against two reduced models of the tip port (run from the repository root with both packages on the path).},label={lst:bt-try}]
import time
import numpy as np
from scipy.linalg import eigh
from fea_engine.material import EI_beam, Material, Section
from fea_engine.geometry import generate_mesh
from fea_engine.solver import FESystem
from fea_engine.damping import RayleighDamping
from fea_engine import elements
from rom_engine.balanced_truncation import BalancedTruncationROM
from rom_engine.krylov import KrylovROM

mesh = generate_mesh(dim=1, L=1.0, n=24)
fes = FESystem(mesh, elements.Beam2DEulerBernoulli())
mat, sec = Material(E=210e9, nu=0.3, rho=7800.0), Section(A=0.01, I=8.33e-6)
fes.assemble_stiffness(EI_beam(mat, sec)); fes.assemble_mass(7800.0 * 0.01)
fes.assemble_damping(RayleighDamping(alpha=2.0, beta=1e-5)); fes.fix_dofs([0], [0, 1])
free = np.asarray(fes.free_dofs)
K, M, C = (a[np.ix_(free, free)] for a in (fes.K, fes.M, fes.C))
B = np.zeros(len(free)); B[len(free) - 2] = 1.0        # collocated tip port

lam, Vf = eigh(K, M)
omegas = np.linspace(0.05, 8, 400) * np.sqrt(lam[0])

# naive: one complex solve per frequency
t0 = time.perf_counter()
H = np.array([B @ np.linalg.solve(-w**2 * M + 1j * w * C + K, B) for w in omegas])
t_full = time.perf_counter() - t0

# reduced: balanced truncation after a 15-mode pre-reduction, and Krylov at s0 = 0
Vp = Vf[:, :15]
bt = BalancedTruncationROM.from_MCK(Vp.T @ M @ Vp, Vp.T @ K @ Vp, Vp.T @ B, Vp.T @ B,
                                    C=Vp.T @ C @ Vp, r=8)
kr = KrylovROM.from_MCK(M, K, B, B, C=C, s0=0.0, k=8)
for name, m in (("balanced truncation", bt), ("Krylov", kr)):
    t0 = time.perf_counter(); Hr = m.frequency_response(omegas); t = time.perf_counter() - t0
    err = np.max(abs(Hr - H)) / np.max(abs(H))
    print(f"{name:20s} 8 states: sweep {1e3*t:5.1f} ms, peak-relative error {err:.1e}")
print(f"full model sweep ({len(free)} unknowns, {len(omegas)} frequencies): {1e3*t_full:.0f} ms")
print(f"a priori bound {bt.h_infinity_error_bound():.1e}, stable: {bt.is_stable()}")
\end{lstlisting}

\begin{lstlisting}[numbers=none,language={},keywordstyle={},commentstyle={},frame=single,aboveskip=2pt]
balanced truncation  8 states: sweep   9.0 ms, peak-relative error 6.4e-06
Krylov               8 states: sweep   7.6 ms, peak-relative error 6.7e-08
full model sweep (48 unknowns, 400 frequencies): 29 ms
a priori bound 3.8e-10, stable: True
\end{lstlisting}

Timings are one run on this machine. With only 48 unknowns the full sweep is already cheap and the reduced sweeps are not much faster, so speed is not the point of this example. The accuracy is: eight states reproduce the sweep to a few parts in $10^{6}$ or better. What is odd is the last line. The bound is $3.8\times10^{-10}$ and the measured error is $6.4\times10^{-6}$, four orders of magnitude larger. That is not a failure of the theorem, because the bound refers to the pre-reduced model and the error here is against the full one; the difference is the price of the modal pre-reduction, taken up in \Cref{sec:bt-raw}.

\begin{table}[ht]
\centering\small
\begin{tabularx}{\linewidth}{@{}l X >{\centering\arraybackslash}p{0.11\linewidth}@{}}
\toprule
Concept & Code & Script calls \\
\midrule
state-space form (descriptor and explicit) of $(\bm M,\bm C,\bm K)$ & \anchor{rom_engine.state_space.to_state_space} & no \\
Gramians, Hankel singular values, balancing transform & \anchor{rom_engine.balanced_truncation.hankel_singular_values} & no \\
balanced truncation, a priori bound \eqref{eq:bt-bound} & \anchor{rom_engine.balanced_truncation.BalancedTruncationROM} & yes \\
singular-perturbation variant (not tested here) & \anchor{rom_engine.balanced_truncation.SingularPerturbationROM} & no \\
one- and two-sided Krylov moment matching & \anchor{rom_engine.krylov.KrylovROM}, \anchor{rom_engine.krylov.arnoldi_basis} & part \\
second-order Arnoldi & \anchor{rom_engine.soar.SOARROM}, \anchor{rom_engine.soar.soar_basis} & yes \\
passivity diagnostic & \anchor{rom_engine.passivity.passivity_margin}, \anchor{rom_engine.passivity.velocity_transfer_function} & yes \\
worked example & \path{rom_engine/examples/gallery/rom_04_balanced_truncation.py} & -- \\
\bottomrule
\end{tabularx}
\caption{Where the concepts of this chapter live in the repository (tag \booktag). The last column says whether this chapter's own checks call, by name, the code listed in the row: yes (every name), part, or no. Code reached only through another routine counts as no, so the column is a lower bound on what this book exercises.}\label{tab:bt-map}
\end{table}

\section{Verification}\label{sec:bt-verif}

All numbers were measured on \machine{}. Errors are measured against the full-order sweep over the band $0.05$--$8\,\omega_1$; two summaries are used, the largest error divided by the peak $|H|$ (``peak-relative'') and the mean of $|H_r-H|/|H|$ over the band (``mean relative''), because the first is dominated by the resonance and the second by the frequencies between resonances.

\subsection{Pre-reduction and the raw model}\label{sec:bt-raw}

The example applies balanced truncation not to the \cbtNState{}-state model but to the $2\times\cbtNPre$-state model obtained by truncating to the lowest \cbtNPre{} undamped modes, and the package documentation says that the raw model is numerically unreliable for the dense Lyapunov solver, because the frequencies span many orders of magnitude. We tested this. On the raw model the order-8 balanced truncation has a peak-relative error of \cbtRawErrMax and a first Hankel singular value of \cbtRawHsvOne (against \cbtPreHsvOne on the pre-reduced model, a difference that should not exist if the two models had the same input--output map to the accuracy of the modal truncation); on the pre-reduced model the same order gives \cbtBtMaxEight. The raw model's reduced system is stable. The pre-reduction itself costs accuracy: the \cbtNPre{}-mode model has a peak-relative error of \cbtPreErrMax against the full sweep, which is the floor for every balanced model of this chapter, and a mean relative error of \cbtPreErrMean.

\subsection{The a priori bound}

On the pre-reduced model the Hankel singular values fall to \cbtHsvRatioTen of the largest at the tenth and \cbtHsvRatioTwenty at the twentieth (\Cref{fig:bt-bt}). \Cref{eq:bt-bound} holds at \cbtBoundHolds{} of the \cbtBoundCases{} orders tested ($r=2,4,\dots,28$), where the actual error is the largest difference to the exact response of the \cbtNPre{}-mode model over the $400$ frequencies. The one exception, $r=\cbtMissOrder$, has an error of \cbtMissErr and a bound of \cbtMissBound, which is \cbtMissRel{} of the largest singular value: rounding level, not a counter-example. Where the bound is meaningful the actual error is between \cbtBoundRatioMin{} and \cbtBoundRatioMax{} of it, so the bound is safe and loose by up to a factor $25$ in this test; at $r=8$ it is \cbtBtBoundEight against an actual error of \cbtBtErrIntEight. A grid maximum is a lower bound on the $H_\infty$ norm, so this checks the bound at the sampled frequencies and does not prove it. The reduced model is stable at all \cbtBtStable{} orders, as the theorem says.

\begin{figure}[ht]
\centering
\includegraphics[width=0.98\linewidth]{figures/ch09_bt.pdf}
\caption{Left: Hankel singular values of the pre-reduced model. Right: the a priori bound \eqref{eq:bt-bound} and the actual grid-maximum error against the retained order. The error is measured against the exact response of the pre-reduced model.}\label{fig:bt-bt}
\end{figure}

\subsection{Accuracy at equal state count}

\Cref{fig:bt-compare} and \Cref{tab:bt-compare} compare five reductions of the same port at the same number of states: modal truncation with $r/2$ undamped modes, balanced truncation, one-sided Krylov at $s_0=0$ on the full model, two-sided Krylov with $r/2$ vectors per side, and SOAR with $r/2$ unknowns.

\begin{figure}[ht]
\centering
\includegraphics[width=0.98\linewidth]{figures/ch09_compare.pdf}
\caption{Error in the band $0.05$--$8\,\omega_1$ against the reduced state count: peak-relative (left) and mean relative (right).}\label{fig:bt-compare}
\end{figure}

\begin{table}[ht]
\centering\small
\begin{tabular}{lrrrr}
\toprule
 & \multicolumn{2}{c}{8 states} & \multicolumn{2}{c}{12 states} \\
\cmidrule(lr){2-3}\cmidrule(lr){4-5}
method & peak-rel. & mean rel. & peak-rel. & mean rel. \\
\midrule
modal truncation & \cbtModalMaxEight & \cbtModalMeanEight & \cbtModalMaxTwelve & \cbtModalMeanTwelve \\
balanced truncation & \cbtBtMaxEight & \cbtBtMeanEight & \cbtBtMaxTwelve & \cbtBtMeanTwelve \\
Krylov, one-sided & \cbtKryOneMaxEight & \cbtKryOneMeanEight & \cbtKryOneMaxTwelve & \cbtKryOneMeanTwelve \\
Krylov, two-sided & \cbtKryTwoMaxEight & \cbtKryTwoMeanEight & \cbtKryTwoMaxTwelve & \cbtKryTwoMeanTwelve \\
SOAR & \cbtSoarMaxEight & \cbtSoarMeanEight & \cbtSoarMaxTwelve & \cbtSoarMeanTwelve \\
\bottomrule
\end{tabular}
\caption{Error against the full-order sweep at equal state count, in the band $0.05$--$8\,\omega_1$.}\label{tab:bt-compare}
\end{table}

Three results stand out. First, one-sided Krylov is the most accurate method in this band at $8$ and $12$ states (SOAR overtakes it in mean relative error from $16$): a mean relative error of \cbtKryOneMeanEight at $8$ states against \cbtBtMeanEight for balanced truncation. It is matched at $s_0=0$ and the band starts near $s_0$; the next subsection shows what that costs. Second, balanced truncation is only modestly better than modal truncation in peak-relative error at $8$ states and about seven times better at $12$, while the mean relative errors are of the same order (\cbtBtMeanEight{} against \cbtModalMeanEight{} at $8$ states). The bound \eqref{eq:bt-bound} is absolute, so it controls the large values near resonance and says little about the small values between resonances; we did not test why the in-between error is larger, and offer as a hypothesis only that the static flexibility of the discarded modes matters there. Third, two-sided Krylov, which is expected to match twice as many moments, is \emph{less} accurate than one-sided at equal state count here ($8$ and $12$ states) and only reaches comparable accuracy at $16$; we saw no benefit of two-sided projection on this port, and the documentation's warning about its fragility is the reason not to prefer it by default.

\subsection{Locality of moment matching}

\Cref{fig:bt-locality} extends the band to $25\,\omega_1$, which contains the third resonance, for models of order $8$ built as before. Inside the band used for the comparison the mean relative error of one-sided Krylov is \cbtNearKryOne; above $12\,\omega_1$ it is \cbtFarKryOne, and balanced truncation shows \cbtNearBt inside and \cbtFarBt above. With only $8$ states neither method represents the third mode; the difference is how the error depends on the distance from the expansion point, rising steeply for Krylov and being roughly flat for the balanced model. The mean relative errors above $12\,\omega_1$ of the two-sided Krylov model and of SOAR are \cbtFarKryTwo and \cbtFarSoar.

\begin{figure}[ht]
\centering
\includegraphics[width=0.62\linewidth]{figures/ch09_locality.pdf}
\caption{Relative error against frequency for models of order $8$ on an extended band; the dotted line marks the upper end of the band of \Cref{tab:bt-compare}.}\label{fig:bt-locality}
\end{figure}

\begin{figure}[ht]
\centering
\includegraphics[width=0.62\linewidth]{figures/ch09_response.pdf}
\caption{Tip frequency response of the full model and of the models of order $8$.}\label{fig:bt-response}
\end{figure}

\subsection{A scale dependence in SOAR}\label{sec:bt-soar}

SOAR is the one method here that needed care. Called on the model in SI units with a unit force, \code{soar_basis(M, K, B, 4, C=C)} returns \cbtSoarRawCols{} vector instead of $4$. The cause is read from the source: a candidate direction is dropped if its norm is below an \emph{absolute} tolerance (\code{tol=1e-12}), and with $\|\bm K^{-1}\bm b\|$ equal to \cbtStaticNorm{} and a recurrence that contracts by about $\|\bm K^{-1}\bm M\|\sim\omega_1^{-2}$ per step, the second and later vectors fall below it. The package reports this through the shape of $\bm V$ and not by an error, so a caller that does not check the shape builds a one-vector model. Rescaling the load to $\|\bm K^{-1}\bm b\|=1$ and time to $\sigma=s/\omega_1$ (which replaces $\bm M,\bm C$ by $\omega_1^2\bm M,\omega_1\bm C$ and leaves the Krylov subspace about $s_0=0$ unchanged) returns \cbtSoarScaledCols{} vectors, and all SOAR results of this chapter use it (\Cref{lst:bt-run}). We did not change the package; a relative tolerance would be the fix, and \Cref{lst:bt-run} shows the work-around.

With the rescaling SOAR converges, but not as fast as one-sided Krylov at the same state count: at $8$ states its mean relative error is \cbtSoarMeanEight against \cbtKryOneMeanEight, at $12$ states \cbtSoarMeanTwelve against \cbtKryOneMeanTwelve. What SOAR gives in exchange is a reduced model that is again second-order with symmetric $\bm M_r,\bm C_r,\bm K_r$, which can be coupled to other components, and the passivity property of the next subsection.

\subsection{Stability and passivity}

\Cref{tab:bt-stab} reports, for $8$, $12$ and $16$ states, whether all poles of the reduced model are in the left half-plane (for the first-order models; the second-order Galerkin models are stable by congruence for damped $\bm K,\bm M,\bm C\succeq0$ and are marked with a dash), and the smallest phase cosine of the mobility over $0.05$--$25\,\omega_1$. All \cbtRowCount{} reduced models have stable poles where this was checked, so we found none of the instability that the package documentation reports for a $k=20$ Krylov model in a published benchmark (\emph{quoted from the documentation, not reproduced here}); and the passivity margin is positive for all of them, with a minimum of about \cbtPassFull for the full model and for every reduced model. That minimum sits at the lowest frequency of the sweep, where all models agree and the mobility of a lightly damped structure is almost purely imaginary, so the table says that no violation was found on this grid, not that the margins are large. Passivity is guaranteed only for the second-order Galerkin models (modal truncation, SOAR); for balanced truncation and Krylov this is evidence from one port and one grid.

\begin{table}[ht]
\centering\small
\begin{tabular}{rlcr}
\toprule
states & method & poles in left half-plane & min $\cos\arg H_v$ \\
\midrule
8 & modal truncation & -- & $4.5\times10^{-4}$ \\
8 & balanced truncation & yes & $4.5\times10^{-4}$ \\
8 & Krylov, one-sided & yes & $4.5\times10^{-4}$ \\
8 & Krylov, two-sided & yes & $4.5\times10^{-4}$ \\
8 & SOAR & -- & $4.5\times10^{-4}$ \\
12 & modal truncation & -- & $4.5\times10^{-4}$ \\
12 & balanced truncation & yes & $4.5\times10^{-4}$ \\
12 & Krylov, one-sided & yes & $4.5\times10^{-4}$ \\
12 & Krylov, two-sided & yes & $4.5\times10^{-4}$ \\
12 & SOAR & -- & $4.5\times10^{-4}$ \\
16 & modal truncation & -- & $4.5\times10^{-4}$ \\
16 & balanced truncation & yes & $4.5\times10^{-4}$ \\
16 & Krylov, one-sided & yes & $4.5\times10^{-4}$ \\
16 & Krylov, two-sided & yes & $4.5\times10^{-4}$ \\
16 & SOAR & -- & $4.5\times10^{-4}$ \\
\bottomrule
\end{tabular}

\caption{Stability and the smallest phase cosine of the mobility ($0.05$--$25\,\omega_1$) for models of $8$, $12$ and $16$ states; a dash marks second-order Galerkin models.}\label{tab:bt-stab}
\end{table}

\section{Limits and failure modes}\label{sec:bt-limits}

\begin{limitsbox}
\begin{itemize}
\item \textbf{One port, one structure.} Everything here is for a $24$-element damped cantilever with a collocated tip port and one band. The ranking of the methods may change for other ports, several inputs or dense modal spectra.
\item \textbf{Balanced truncation is dense.} The Lyapunov solves are cubic in the number of states, the raw \cbtNState{}-state model gave an unreliable result (\Cref{sec:bt-raw}), and the modal pre-reduction used instead sets the accuracy floor. No large-scale Gramian solver is tested.
\item \textbf{The bound is absolute and loose}, by up to a factor of about $25$ at the sampled frequencies, and a grid check is not a proof.
\item \textbf{Moment matching has no error control away from $s_0$} (\Cref{fig:bt-locality}), and two-sided projection gave no benefit here.
\item \textbf{SOAR needs a rescaled problem}, because of an absolute tolerance that reports through the returned shape; TOAR is not implemented. Passivity is a theorem only for second-order Galerkin models; for the others it is a grid check.
\end{itemize}
Not tested: the singular-perturbation, frequency-weighted and Hankel-norm variants, several-input ports, models above \cbtNState{} states, and timing beyond the single run of \Cref{lst:bt-try}.
\end{limitsbox}

\section{Summary}\label{sec:bt-summary}

We started with a loop over frequencies and asked what a component really needs from the beam, which is the map from one force to one displacement. Balanced truncation reduces that map with a stability guarantee and a bound given before any response is computed; the bound held at every meaningful order and was loose, and it costs a dense Lyapunov solve and, here, a modal pre-reduction. Krylov moment matching skips the Lyapunov solve and was the most accurate method in the band around its expansion point, and the least accurate away from it. SOAR stays in second-order form, which keeps passivity by congruence, but its package implementation needs the problem rescaled before it returns a full basis.

The practical choice follows from what is known. A narrow, known band points to Krylov or SOAR; a need for a global guarantee points to balanced truncation; a reduced model that will be coupled to other components points to SOAR or modal truncation. Two-sided Krylov gave no benefit on this port. The next chapter keeps the frequency domain but chooses the sampling points by a greedy search with a certified error bound.

\section*{Exercises}\label{sec:bt-exercises}
\begin{enumerate}
\item \emph{(Implementation)} Repeat \Cref{tab:bt-compare} with the expansion point moved to $s_0=4\omega_1$ (a real shift). Where does the one-sided Krylov error now sit, and does the model stay stable?
\item \emph{(Implementation)} Apply balanced truncation to the model pre-reduced to $10$ and to $25$ modes. How does the accuracy floor and the order needed for a given error change?
\item \emph{(Interpretation)} Use a non-collocated port (force at the tip, displacement at midspan). Does the passivity argument of \Cref{sec:bt-math} still apply, and what does \code{passivity_margin} report?
\item \emph{(Implementation)} Replace the absolute tolerance in \code{soar_basis} by a tolerance relative to the norm of the first vector, and check that the unscaled call of \Cref{sec:bt-verif} returns four vectors. (Test file: \texttt{rom\_engine/tests/test\_soar.py}.)
\item \emph{(Implementation)} Find the order at which \eqref{eq:bt-bound} would first guarantee a peak-relative error below $10^{-6}$ for the pre-reduced model, and compare with the order the measured error needs.
\item \emph{(Derivation)} For $\bm E\dot{\bm x}=\bm A\bm x+\bm B u$, show that if $\bm V$ spans the Krylov subspace $\mathcal K_r\big((\bm A-s_0\bm E)^{-1}\bm E,(\bm A-s_0\bm E)^{-1}\bm B\big)$ and $\bm W=\bm V$, the first $r$ Taylor coefficients of the reduced transfer function about $s_0$ equal those of the full one.
\end{enumerate}

\chapterrefs
